%! Adding \& Subtracting Rational Expressions — Unit 4, Lesson 4.3 — Guided Notes
\documentclass[10pt]{article}
\usepackage{algebra2-article}
\usepackage{algebra2-boxes}
\newcommand{\TallMath}[1]{$\displaystyle #1\rule[-1.4em]{0pt}{3.2em}$}
\newcommand{\lgrid}[4]{%
  \draw[step=1, linegray!40, very thin] (#1,#3) grid (#2,#4);
  \draw[->, semithick, charcoal] (#1,0)--(#2,0) node[below right, font=\scriptsize]{$x$};
  \draw[->, semithick, charcoal] (0,#3)--(0,#4) node[left, font=\scriptsize]{$y$};}

\begin{document}

\pageheader{Unit 4, Lesson 4.3}{Guided Notes}
\namedateperiod

\vspace{0.08in}

\begin{objectivebox}
\textbf{By the end of this lesson, I will be able to\dots}
\begin{itemize}\setlength{\itemsep}{2pt}
  \item add and subtract rational expressions with like denominators, putting the second numerator in \emph{parentheses} whenever I subtract;
  \item build the \textbf{least common denominator} from the \emph{factored} denominators, rewrite each fraction to match it, combine, and simplify;
  \item recognize a \textbf{complex fraction} and simplify it as a quotient of two simple fractions --- and state every restriction it carries.
\end{itemize}
\end{objectivebox}

\vspace{0.05in}

\begin{vocabbox}
\small
Fill in each term as we build it together.
\par\vspace{2pt}   % \par is required: \termblanklong's \noindent is a no-op mid-paragraph
\termblanklong{Least common denominator (LCD)}
\termblanklong{Building factor}
\termblanklong{Complex algebraic fraction}
\termblanklong{Main fraction bar}
\end{vocabbox}

\begin{hookbox}
Two students subtracted the same expression and got different answers. Let
\[
  D(x)=\frac{4x-1}{x+3}-\frac{2x-7}{x+3} .
\]
\begin{center}
\renewcommand{\arraystretch}{1.35}
\begin{tabular}{|l|c|c|c|}
\hline
 & $x=0$ & $x=1$ & $x=-3$ \\ \hline
$D(x)$ \; (the original) & \blank{1.2cm} & \blank{1.2cm} & \blank{1.2cm} \\ \hline
Dev: \; $\dfrac{2x-8}{x+3}$ & \blank{1.2cm} & \blank{1.2cm} & \blank{1.2cm} \\ \hline
Elin: \; $2$ & \blank{1.2cm} & \blank{1.2cm} & \blank{1.2cm} \\ \hline
\end{tabular}
\end{center}
\begin{itemize}\setlength{\itemsep}{1pt}
  \item Dev wrote $\dfrac{4x-1-2x-7}{x+3}$. Which \emph{one} keystroke is missing, and what did losing
        it do to the sign of the $7$? \writeline
  \item Elin is right about the value --- but look at the last column. What must she add to her answer
        before it is finished? \writeline
\end{itemize}
Two ideas run the whole period. \textbf{A minus sign in front of a fraction belongs to that fraction's
entire numerator.} And once again: even an answer as plain as ``$2$'' still carries the restrictions of
the problem it came from.
\end{hookbox}

\begin{notesbox}{1. Like Denominators: Add the Numerators, Keep the Denominator}
When the denominators already match, there is nothing to build:
\[
  \frac{a}{c}+\frac{b}{c}=\frac{a+b}{c}, \qquad
  \frac{a}{c}-\frac{b}{c}=\frac{a-b}{c}, \qquad c\neq0 .
\]

\vspace{3pt}
\textbf{You never add the denominators.} Check it on numbers you already know:
$\dfrac12+\dfrac12=$ \blank{0.8cm}, \; but adding denominators would give
$\dfrac{2}{4}=$ \blank{0.8cm}. Adding denominators would make two halves into
\blank{2.0cm}, which is nonsense.

\vspace{4pt}
\textbf{Worked.} \; $\dfrac{7}{x^2-4}+\dfrac{x}{x^2-4}$
\begin{itemize}[leftmargin=1.5em, itemsep=3pt]
  \item Factor the denominator \emph{first}, so you can see the restrictions: $x^2-4=$ \blank{2.4cm},
        so $x\neq$ \blank{1.8cm}.
  \item Add the numerators, keep the denominator: \blank{3.0cm}
  \item \textbf{Last step, always:} factor the new numerator and check whether anything divides out.
        Here $x+7$ is prime, so we are done.
\end{itemize}
\end{notesbox}

\begin{notesbox}{2. Subtracting: the Minus Sign Owns the Whole Numerator}
\begin{center}
\fbox{\parbox{0.90\linewidth}{\centering\vspace{2pt}
\textbf{Put the second numerator in parentheses before you subtract.} Then distribute the minus sign
--- it changes the sign of \emph{every} term inside.\vspace{2pt}}}
\end{center}

\vspace{4pt}
\textbf{Worked.} \; $\dfrac{6x+5}{x^2-4}-\dfrac{2x+13}{x^2-4}$
\begin{itemize}[leftmargin=1.5em, itemsep=3pt]
  \item Write the subtraction with parentheses: $\dfrac{(6x+5)-(\rule{2.0cm}{0.4pt})}{x^2-4}$
  \item Distribute and collect: numerator $=$ \blank{3.0cm} $=$ \blank{2.4cm} (factored).
  \item Divide out, then attach the restrictions from the \emph{original} denominator:
        \blank{2.0cm}, \; $x\neq$ \blank{1.8cm}.
  \item Check it at $x=0$: the original gives \blank{1.2cm} and your answer gives \blank{1.2cm}.
\end{itemize}

\vspace{3pt}
Notice which restriction survived. The factor $(x-2)$ divided out, so $x=2$ is a
\blank{1.4cm} in the graph; $(x+2)$ stayed, so $x=-2$ is a \blank{1.4cm}. That sort is Lesson 4.0's,
and it never stops applying.
\end{notesbox}

\begin{notesbox}{3. Unlike Denominators: Build the LCD}
You cannot combine pieces that are cut differently. The \textbf{least common denominator} is the
smallest expression every denominator divides into --- and you build it out of \emph{factors}, never by
multiplying the denominators together.

\vspace{4pt}
\textbf{The six steps.}
\begin{enumerate}[leftmargin=1.9em, itemsep=2pt]
  \item \textbf{Factor} every denominator completely.
  \item \textbf{Build the LCD:} take every \emph{distinct} factor, each to the \emph{highest} power it reaches in any one denominator.
  \item \textbf{List the restrictions} --- set each factor of the LCD equal to zero.
  \item \textbf{Rewrite} each fraction: multiply it by the \textbf{building factor} $\frac{k}{k}$ that supplies whatever its denominator is missing.
  \item \textbf{Combine} the numerators over the LCD (parentheses on every subtraction) and simplify.
  \item \textbf{Factor the numerator} and divide out if you can. Attach the restrictions.
\end{enumerate}

\vspace{4pt}
\textbf{Monomial denominators.} \; $\dfrac{5}{6x^2}+\dfrac{7}{4x}$
\begin{itemize}[leftmargin=1.5em, itemsep=3pt]
  \item $6x^2=2\cdot3\cdot x^2$ and $4x=2^2\cdot x$, so LCD $=$ \blank{1.6cm}.
  \item Building factors: the first needs \blank{1.0cm}, the second needs \blank{1.2cm}.
  \item Result: \blank{2.6cm}, \; $x\neq$ \blank{1.0cm}.
\end{itemize}

\vspace{5pt}
\textbf{The anchor.} \; $\dfrac{3}{x-3}-\dfrac{18}{x^2-9}$
\begin{itemize}[leftmargin=1.5em, itemsep=3pt]
  \item Factor: $x^2-9=$ \blank{2.4cm}, so LCD $=$ \blank{2.8cm} and $x\neq$ \blank{1.8cm}.
  \item The first fraction is missing the factor \blank{1.4cm}; multiply it by that over itself:
        $\dfrac{3}{x-3}\cdot\dfrac{\rule{1.2cm}{0.4pt}}{\rule{1.2cm}{0.4pt}}=
        \dfrac{\rule{1.8cm}{0.4pt}}{(x-3)(x+3)}$
  \item Combine and simplify the numerator: \blank{2.4cm} $=$ \blank{2.4cm} (factored).
  \item Divide out. Final answer: \blank{2.0cm}, \; $x\neq$ \blank{1.8cm}.
\end{itemize}
Where have you seen that answer before today? \writeline

\vspace{5pt}
\textbf{Two trinomial denominators.} \; $\dfrac{2}{x^2+5x+6}+\dfrac{3}{x^2-9}$
\begin{itemize}[leftmargin=1.5em, itemsep=3pt]
  \item Factored: \blank{2.6cm} and \blank{2.6cm}. The shared factor is \blank{1.4cm}, so the LCD uses
        it \emph{once}: LCD $=$ \blank{4.0cm}.
  \item Restrictions: $x\neq$ \blank{2.6cm}
  \item Combine: $\dfrac{2(\rule{1.2cm}{0.4pt})+3(\rule{1.2cm}{0.4pt})}{\text{LCD}}=$ \blank{3.4cm}
\end{itemize}

\vspace{4pt}
\begin{center}
\fbox{\parbox{0.92\linewidth}{\centering\vspace{2pt}
\textbf{The LCD \emph{is} the restriction list.} Every factor of every original denominator is in
there, so setting the LCD's factors to zero gives you every excluded value --- and it costs you no
extra work.\vspace{2pt}}}
\end{center}
\vspace{2pt}
One warning attached to that shortcut: build the LCD from the \emph{original} factored denominators.
Reading anything off your final answer is the error from Lesson 4.1, and it is still wrong today.

\vspace{5pt}
\textbf{Opposite binomials in the denominators} (Lesson 4.1, still in force). \;
$\dfrac{4}{x-5}+\dfrac{2}{5-x}$
\begin{itemize}[leftmargin=1.5em, itemsep=3pt]
  \item These denominators are not different --- they are opposites: $5-x=$ \blank{1.2cm} $(x-5)$.
  \item Rewrite the second fraction as $\dfrac{\rule{1.0cm}{0.4pt}}{x-5}$. Now the denominators match.
  \item Answer: \blank{2.0cm}, \; $x\neq$ \blank{1.0cm}. \; (The LCD was never $(x-5)(5-x)$.)
\end{itemize}
\end{notesbox}

\begin{notesbox}{4. Complex Fractions: a Division in Disguise}
A \textbf{complex fraction} has fractions inside its numerator, its denominator, or both. The long
\textbf{main fraction bar} is not decoration --- it is a division sign.
\[
  \frac{\;\dfrac{1}{x}+\dfrac{1}{2}\;}{\;\dfrac{1}{x}-\dfrac{1}{4}\;}
  \qquad\text{means}\qquad
  \left(\frac{1}{x}+\frac{1}{2}\right)\div\left(\frac{1}{x}-\frac{1}{4}\right)
\]

\vspace{3pt}
\textbf{Method 1 --- combine, then divide.} Turn the top into one simple fraction, turn the bottom into
one simple fraction, then keep--change--flip (Lesson 4.2).
\begin{itemize}[leftmargin=1.5em, itemsep=3pt]
  \item Top: $\dfrac{1}{x}+\dfrac{1}{2}=$ \blank{2.0cm} \qquad Bottom: $\dfrac{1}{x}-\dfrac{1}{4}=$
        \blank{2.0cm}
  \item Divide: \; \blank{2.0cm} $\cdot$ \blank{2.0cm} $=$ \blank{2.4cm}
\end{itemize}

\vspace{3pt}
\textbf{Method 2 --- clear all the little fractions at once.} Multiply the top \emph{and} the bottom by
the LCD of every small fraction inside. Here that LCD is \blank{1.2cm}.
\begin{itemize}[leftmargin=1.5em, itemsep=3pt]
  \item Top becomes \blank{1.6cm}; \; bottom becomes \blank{1.6cm}; \; so the answer is
        \blank{2.2cm} --- the same one.
\end{itemize}

\vspace{4pt}
\textbf{Restrictions: two levels, and the second one hides.}
\begin{center}
\renewcommand{\arraystretch}{1.5}
\begin{tabularx}{\linewidth}{@{}c l X@{}}
\hline
\textbf{\#} & \textbf{What must be nonzero} & \textbf{Why --- and where it hides} \\
\hline
1 & every \emph{inner} denominator & each little fraction has to exist. Visible: just read them. Here $x\neq$ \blank{0.8cm}. \\
2 & the \emph{whole bottom} & the main bar is a $\div$, and you cannot divide by zero. This is Lesson 4.2's \textbf{source 3} in a new costume: set the bottom's combined \emph{numerator} to zero. Here $4-x=0$ gives $x\neq$ \blank{0.8cm}. \\
\hline
\end{tabularx}
\end{center}
Complete restriction list for the anchor: $x\neq$ \blank{2.0cm}. Check the whole thing at $x=2$: the
original gives \blank{1.0cm} and your answer gives \blank{1.0cm}.

\vspace{5pt}
\textbf{One more, with a binomial underneath.} \;
$\dfrac{\;\dfrac{1}{x}-\dfrac{1}{3}\;}{\;x-3\;}$
\begin{itemize}[leftmargin=1.5em, itemsep=3pt]
  \item Top: \blank{1.8cm}. \; The bottom is already one piece, so divide by it:
        \blank{1.8cm} $\cdot\;\dfrac{1}{x-3}$
  \item Watch the \emph{opposite binomials}: $3-x=$ \blank{1.0cm} $(x-3)$, so the answer is
        \blank{1.8cm}.
  \item Restrictions: $x\neq$ \blank{1.6cm} --- one from the inner denominator, one from the main bar.
\end{itemize}
\end{notesbox}

\begin{notesbox}{5. The Picture: One Subtraction, One Hole and One Wall}
The anchor from Section 3 was $\;g(x)=\dfrac{3}{x-3}-\dfrac{18}{x^2-9}=\dfrac{3}{x+3}$, with
$x\neq3,-3$.
\begin{center}
\begin{tikzpicture}[scale=0.42]
  \lgrid{-8}{5}{-5}{5}
  \draw[dashed, redacc, thick] (-3,-5)--(-3,5);
  \node[redacc, font=\scriptsize, above left] at (-3.1,3.4) {$x=-3$};
  \begin{scope}
    \clip (-8,-5) rectangle (5,5);
    \draw[very thick, forest, domain=-8:-3.62, samples=100, smooth]
      plot (\x,{3/((\x)+3)});
    \draw[very thick, forest, domain=-2.38:5, samples=100, smooth]
      plot (\x,{3/((\x)+3)});
  \end{scope}
  \draw[forest, very thick, fill=white] (3,0.5) circle (4pt);
  \node[charcoal, font=\scriptsize, above right] at (3.1,0.7) {hole};
\end{tikzpicture}
\end{center}
\begin{enumerate}[leftmargin=1.7em, itemsep=3pt]
  \item Both excluded values came from the \emph{original} denominators. Write those two denominators
        in factored form: \blank{2.0cm} and \blank{2.6cm}. \; Which excluded value appears in
        \emph{both} of them? \blank{1.2cm}
  \item $x=3$ is a \blank{1.4cm} because its factor \blank{2.4cm}; \; $x=-3$ is a
        \blank{1.4cm} because its factor \blank{2.4cm}.
  \item The hole's height comes from the \emph{simplified} form: $\dfrac{3}{3+3}=$ \blank{1.0cm}, so
        the hole sits at \blank{1.6cm}.
  \item Your final answer $\dfrac{3}{x+3}$ shows you only one of the two restrictions. Explain, in one
        sentence, why the other one still belongs to it. \writeline
\end{enumerate}
\end{notesbox}

\begin{practicebox}
Simplify completely and state \emph{all} restrictions. Show the LCD you used.

\vspace{3pt}
\begin{enumerate}[leftmargin=1.7em, itemsep=8pt]
  \item $\dfrac{x}{x+4}+\dfrac{4}{x+4}$ \hfill Simplest form: \blank{2.0cm} \quad
        $x\neq$ \blank{1.4cm} \\[3pt]
        Your answer is a plain number. What is left of the fraction it came from? \blank{4.6cm}
  \item $\dfrac{5}{x}-\dfrac{2}{x+1}$ \hfill LCD: \blank{2.4cm} \\[4pt]
        Simplest form: \blank{2.6cm} \quad $x\neq$ \blank{2.0cm}
  \item $\dfrac{\;\dfrac{2}{x}+3\;}{\;\dfrac{2}{x}-3\;}$ \hfill Simplest form: \blank{2.4cm} \quad
        $x\neq$ \blank{2.2cm} \\[3pt]
        Which of your two restrictions came from the \emph{main} fraction bar? \blank{2.0cm}
\end{enumerate}
\end{practicebox}

\end{document}
